% Technical supplement for the campus LLM policy agent.
% The prompt wording below follows covasim/llm_policy.py. Whitespace is
% normalized only for typesetting.

\section{Technical Supplement: LLM Policy Agent}
\label{app:llm-policy-agent}

\subsection{Role of the language model}

The language model is used as a constrained, high-level policy selector inside
the campus epidemic simulator.  It does not generate epidemiological parameters
or directly manipulate the contact network.  At each decision time, the agent
receives an aggregate state summary and selects exactly one policy from the
ordered catalog
\(
\mathcal{P}=\{P0_{\mathrm{normal}},P1_{\mathrm{guarded}},
P2_{\mathrm{moderate}},P3_{\mathrm{strong}},P4_{\mathrm{emergency}}\}
\).
The simulator, rather than the LLM, maps the selected identifier to fixed
testing, screening, contact-tracing, teaching, and contact-reduction parameters.
Thus, the action space is finite and all low-level intervention parameters are
specified before inference.

In the paper configuration, decisions are made every seven simulation days.
The initial policy is \texttt{P0\_normal}, and the first model decision is
deferred until day 7, so that one complete observation window is available.
All features are computed using simulation days preceding the decision day;
the current decision day is not included in the seven-day window.

\subsection{Fixed policy action space}

Table~\ref{tab:policy-catalog-brief} reports the policy information exposed to
the LLM.  The detailed numerical intervention profiles are deliberately omitted
from the prompt.  They remain fixed in the simulator and cannot be edited by the
model.

\begin{table}[h]
\centering
\small
\begin{tabular}{clccl}
\hline
Rank & Policy ID & Intensity & Online & Prompt-level description \\
\hline
0 & \texttt{P0\_normal} & 0 & No & Normal teaching and baseline controls \\
1 & \texttt{P1\_guarded} & 1 & No & Stronger screening/tracing; mild restrictions \\
2 & \texttt{P2\_moderate} & 2 & No & Higher screening; reduced in-person intensity \\
3 & \texttt{P3\_strong} & 3 & No & High-cost, hybrid-like strong control \\
4 & \texttt{P4\_emergency} & 4 & Yes & Maximum restrictions and online teaching \\
\hline
\end{tabular}
\caption{Compact policy catalog supplied to the LLM.  The implementation uses
longer natural-language summaries with the same meaning.}
\label{tab:policy-catalog-brief}
\end{table}

For exact reproducibility, the five \texttt{summary} strings serialized in
\texttt{policy\_catalog\_brief} are:

\begin{itemize}
  \item \texttt{P0\_normal}: ``Lowest-cost posture with normal teaching,
  baseline symptom testing, low-frequency targeted screening, and routine
  tracing.''
  \item \texttt{P1\_guarded}: ``Low-cost guarded posture with stronger
  screening and tracing plus mild public-space restrictions while keeping
  in-person teaching normal.''
  \item \texttt{P2\_moderate}: ``Medium-cost posture with higher screening
  cadence, strong tracing, and reduced classroom, lab, and office intensity
  while staying mostly in person.''
  \item \texttt{P3\_strong}: ``High-cost strong-control posture with
  hybrid-like contact reduction, elevated screening, and strict campus
  restrictions to counter clearly worsening spread.''
  \item \texttt{P4\_emergency}: ``Highest-cost emergency posture with online
  teaching, campus-wide high-frequency screening, and maximum restrictions for
  severe uncontrolled spread.''
\end{itemize}

\subsection{System prompt}

The following is the complete system prompt used by the full agent.  Wording is
reproduced exactly; whitespace between sentences is normalized for typesetting.
For ablations without expert knowledge, the final paragraph beginning with
``Use these short systems rules'' is omitted.

\begin{quote}\small
You are a campus epidemic prevention and control expert. Your task is to choose
exactly one next-period policy based on the aggregated campus state. You may
only choose one of these \texttt{policy\_id} values:
\texttt{P0\_normal}, \texttt{P1\_guarded}, \texttt{P2\_moderate},
\texttt{P3\_strong}, \texttt{P4\_emergency}. Do not create new policies. Do not
modify policy parameters. Do not output any low-level transmission parameters.
Do not give real-world medical advice. You are only selecting the next
simulation policy.

Output must be strict JSON containing exactly three fields:
\texttt{policy\_id}, \texttt{risk\_level}, \texttt{reason}.
\texttt{policy\_id} must be one of the five fixed policies.
\texttt{risk\_level} must be one of \texttt{low}, \texttt{medium},
\texttt{high}, \texttt{critical}. \texttt{reason} must be a short Chinese
sentence within 40 characters.

Key data fields in the state you receive:
\texttt{growth\_ratio\_7d\_vs\_prev\_7d} is the ratio of new infections in the
past 7 days versus the previous 7 days. Above 1.0 means the epidemic is growing
week-over-week; below 1.0 means it is shrinking.
\texttt{new\_infections\_per\_1000\_last\_7d} is the weekly new infections per
1000 population. \texttt{new\_infections\_trend} is the direction label
(\texttt{up}/\texttt{down}/\texttt{flat}) comparing the past 7 days to the
previous 7 days. \texttt{estimated\_risk\_trend} is the overall risk direction
(\texttt{improving}/\texttt{stable}/\texttt{worsening}).
\texttt{layer\_transmission\_trends} are per-layer trend labels showing which
contact layers are rising or falling. \texttt{recent\_policy\_memory} is the
history of recent policy changes and their observed effects.
\texttt{short\_term\_trajectory\_summary} describes recent policy evolution,
infection trend, growth ratio, and the observed feedback pattern after the
latest policy adjustment. \texttt{similar\_scenarios\_summary} describes
historical situations with similar epidemic states, grouped by the policy that
was executed, showing the observed outcomes for each policy choice.

Use these short systems rules when deciding: Epidemic indicators often lag
policy changes, so do not downgrade too quickly after a single-day decline.
Base the immediate decision primarily on the recent 7-day versus previous 7-day
trends, and use recent policy memory to judge whether the trend is sustained and
whether the current policy is working across multiple periods. When the epidemic
is clearly growing (rising cases week over week), prioritize suppressing the
contact layers whose transmission trends are worsening the most. Dormitory,
classroom, and community transmission can spill over into each other at the
system level, so compare both current counts and trend direction. If infections,
key layer transmission, and policy effects are all improving, prefer holding the
current policy level to confirm the trend is sustained. Only consider
downgrading when the epidemic shows sustained, unambiguous improvement across
multiple consecutive decision cycles---a single week of decline is not enough
to confirm the trend is durable. \texttt{P3\_strong} and
\texttt{P4\_emergency} have high operational and teaching costs, so reserve them
for clearly worsening transmission under the current policy. Do not escalate
only because one layer is still the largest source; check whether that layer is
rising or falling. Testing, isolation, and online teaching reduce transmission
but also increase operational burden and disruption on the campus. Frequent
policy up-down switching reduces stability and should be avoided unless risk is
clearly worsening.
\end{quote}

\subsection{User message and state representation}

The user message is a UTF-8 JSON serialization of the following template:

\begin{verbatim}
{
  "task": "Choose exactly one next-period policy_id from the predefined
           campus policy catalog.",
  "state": { ... aggregate state fields ... },
  "policy_catalog_brief": [ ... five fixed policy summaries ... ],
  "output_schema": {
    "policy_id": "P0_normal | P1_guarded | P2_moderate |
                  P3_strong | P4_emergency",
    "risk_level": "low | medium | high | critical",
    "reason": "short Chinese reason, <=40 chars"
  }
}
\end{verbatim}

The aggregate state contains the following feature groups:

\begin{itemize}
  \item \textbf{Decision context:} decision day/date, current policy, duration
  under the current policy, current restriction intensity, and whether online
  teaching is active.
  \item \textbf{Seven-day epidemic summary:} sums and means of new infections
  and diagnoses, their categorical trends, test count, a diagnosis-to-test
  positivity proxy, mean quarantine occupancy, and estimated risk trend.
  \item \textbf{Normalized severity and dynamics:} new infections per 1,000,
  diagnosis ascertainment proxy, the ratio of infections in the latest seven
  days to the preceding seven days, and real-time effective reproduction number
  summaries (current value, mean, trend, and source).
  \item \textbf{Contact-layer evidence:} seven-day transmission counts, the
  largest contributing layer, and current-versus-previous-window counts, deltas,
  ratios, and trend labels for classroom, dormitory, office/laboratory, and
  community layers.
  \item \textbf{Intervention burden:} tests and mean quarantined individuals,
  both as totals and per 1,000 population, plus teaching restriction intensity.
  \item \textbf{Recent decision memory:} previous state summaries, proposed and
  executed policies, stated risk, and observed post-decision changes in
  infections, testing, quarantine, and teaching disruption.  The standard
  full-agent configuration retains the most recent four weekly decisions.
  \item \textbf{Memory summaries:} a recency-weighted policy-effect summary and
  a policy-volatility summary.  When dual-branch memory is enabled, a compressed
  short-term trajectory and retrieved long-term similar scenarios are included
  as natural-language summaries.  Their full structured records and the
  internal population-size field are excluded from the LLM message to control
  token usage, but are retained in the audit log.
\end{itemize}

The layer keys used internally are \texttt{d} (dormitory), \texttt{clu} and
\texttt{clg} (classroom/course layers), \texttt{of} and \texttt{ta}
(office/laboratory-related layers), and \texttt{cm} (community/common-space
layer).

\subsection{Required output and validation}

The only accepted model output is a JSON object with exactly three keys:

\begin{verbatim}
{
  "policy_id": "P2_moderate",
  "risk_level": "medium",
  "reason": "传播仍在上升，建议加强防控"
}
\end{verbatim}

The parser rejects missing keys, additional keys, non-object JSON, an unknown
policy identifier, or a risk label outside the four-value vocabulary.  The
reason must be a string; leading and trailing whitespace is removed and text
longer than 40 characters is truncated to 40 characters.  The API additionally
requests JSON-object response formatting.  If the response contains surrounding
text, the parser attempts to extract one JSON object; it also performs limited
brace/bracket balancing for truncated responses.  Empty, filtered, or otherwise
invalid responses are retried up to three times with exponential backoff.  If
all attempts fail, the current policy is retained as the safe fallback and the
error is recorded.

\subsection{Decision post-processing}

The model output is a proposal, not an unconditionally executed action.  The
following deterministic processing is applied after schema validation:

\begin{enumerate}
  \item \textbf{Transition constraint.} Escalation to a stricter policy is
  allowed immediately, whereas relaxation is limited to one rank per weekly
  decision.  Formally, for current rank \(r_t\) and proposed rank \(\hat r\),
  \(r'=\hat r\) if \(\hat r>r_t\), and
  \(r'=\max(\hat r,r_t-1)\) otherwise.
  \item \textbf{Minimum-hold guardrail.} A proposal reversing the most recent
  nonzero policy-change direction is blocked until the current policy has been
  held for two decision intervals (14 days in the full configuration), except
  when the reported risk is \texttt{critical}.
  \item \textbf{Burden guardrail.} A non-critical escalation is blocked when
  operational burden is already high but epidemic evidence is not worsening.
  High burden is triggered by at least 20 mean quarantined persons per 1,000,
  at least 80 tests per 1,000 in seven days, or restriction intensity at least
  3. Worsening evidence overrides this block.
  \item \textbf{Low-prevalence online-teaching guardrail.} An escalation to an
  online-teaching policy is capped at \texttt{P2\_moderate} (or the current
  stricter non-online policy) when weekly incidence is below 10 per 1,000, fewer
  than 500 infections occurred in the window, and the growth ratio is not above
  1.5.
\end{enumerate}

For the optional three-member ensemble, the policy ID is selected by majority
vote.  Ties favor the stricter policy rank.  If multiple responses support the
winning policy, the response with the highest risk level supplies the reported
risk and reason.  Invalid member responses are ignored unless all members fail.

\subsection{Inference and reproducibility settings}

The implementation uses an OpenAI-compatible Chat Completions interface.  For
the experiments in \texttt{t17\_campus\_llm\_thesis\_experiments.py}, the
command-line defaults are temperature 0.2 and maximum output length 4,096
tokens; JSON-object response formatting and model-side thinking are enabled.
The current repository defaults point to model
\texttt{deepseek-v4-flash} at \texttt{https://api.deepseek.com}.  The Python
client class and environment-variable names retain the historical
``Qwen/DashScope'' naming, but these names do not describe the current default
provider.  Each JSONL audit record stores the model name, endpoint, temperature,
token limit, finish reason, token usage, raw response, parsed decision, complete
state snapshot, user payload, exact system prompt, and SHA-256 hash of the system
prompt.

Importantly, the experiment script calls the external model only when the
\texttt{--use-api} flag is supplied.  Without this flag, LLM-labeled strategies
use a deterministic surrogate decision function for smoke testing and simulator
validation; such runs must not be reported as real-LLM evaluations.

\subsection{Prompt and memory ablations}

The code supports controlled ablations that preserve the same action space and
output schema.  Disabling expert knowledge removes only the nine systems rules
from the system prompt.  Disabling recent memory sets the decision-history
window to zero.  The dual-branch memory study independently enables (i) a
compressed within-episode short-term trajectory over four entries and (ii)
long-term retrieval of up to six similar cases, with at most two cases per
executed-policy group.  Similarity combines state, contact-layer, temporal-trend,
and text components with weights 0.30, 0.25, 0.20, and 0.25, respectively.  The
full single-agent policy used in the core comparison enables expert rules,
four-week recent memory, the 14-day reversal hold, and both burden guardrails.
